\documentclass[11pt,a4paper]{article}
\usepackage[utf8]{inputenc}
\usepackage[T1]{fontenc}
\usepackage{amsfonts}
\usepackage[french]{babel}
\frenchbsetup{StandardLists=true}
\usepackage{enumitem}
\usepackage{amssymb}
\usepackage{amsmath}
\usepackage[left=2cm,right=2cm,top=2cm,bottom=2cm]{geometry}
\usepackage{multirow}
\usepackage[dvipsnames]{xcolor}

\usepackage[T1]{fontenc}
\usepackage{fouriernc}
\usepackage{graphicx}

\usepackage{setspace}
\singlespacing

\usepackage{tcolorbox}
\usepackage{multicol}
\usepackage{caption}
\usepackage{fancyhdr}
\pagestyle{fancy}
\renewcommand\headrulewidth{1pt}
\fancyhead[L]{Lycée Denis Diderot }
\fancyhead[R]{Classe de 3PM}
\setcounter{page}{1}\fancyfoot[C]{page \thepage}
\fancyfoot[R]{Année 2025-2026}
\fancyfoot[L]{Alexis RUIZ}
\usepackage{multirow,tabularx,booktabs}

\usepackage{pifont}
\usepackage{chemmacros}
\chemsetup{modules=all}
\setlength{\columnseprule}{.25mm}
\setlength{\parindent}{0pt}

\renewcommand{\arraystretch}{2}

\frenchbsetup{StandardLists=true}

\newcommand{\cadreblanc}[1]{%
	\medskip\framebox[\linewidth][c]{%
		\rule{0mm}{#1mm}}\par
	\medskip}



\begin{document}
	
	\begin{center}
		\LARGE{CONTRÔLE - Chapitre 4 - Les ions}
		
		\vspace{3mm}
		
		\large{Durée : 50 minutes}
		
	\end{center}
	
	\normalsize
	
	\hrule\
	\\[0,2cm]
	
	\textbf{Nom :} \dotfill \hspace{2cm} \textbf{Prénom :} \dotfill
	
	\vspace{3mm}
	
	\begin{tcolorbox}[colback=blue!5!white,colframe=blue!75!black]
		{
			\textbf{Consignes :}
			\begin{itemize}
				\item Répondre sur le sujet
				\item Soigner la présentation et l'orthographe
				\item Justifier les réponses quand demandé
				\item Utiliser un stylo bleu ou noir
			\end{itemize}
		}
	\end{tcolorbox}
	
	\vspace{5mm}
	
	\hrule
	\vspace{3mm}
	
	\textbf{Exercice 1 : Connaissances de cours} \hfill \textbf{/6 points}
	
	\vspace{3mm}
	
	\textbf{1.} Qu'est-ce qu'un ion ? \hfill \textit{(1 pt)}
	
	\cadreblanc{15}
	
	\textbf{2.} Quelle est la différence entre un cation et un anion ? \hfill \textit{(2 pts)}
	
	\cadreblanc{20}
	
	\textbf{3.} Comment un atome se transforme-t-il en ion ? \hfill \textit{(1 pt)}
	
	\cadreblanc{15}
	
	\textbf{4.} Citer deux ions que l'on peut identifier par des tests chimiques. \hfill \textit{(1 pt)}
	
	\cadreblanc{12}
	
	\textbf{5.} Vrai ou Faux : Une solution ionique est toujours électriquement neutre. \hfill \textit{(1 pt)}
	
	\cadreblanc{10}
	
	\vfill
	
	\hrule
	\vspace{3mm}
	
	\newpage
	
	
	\textbf{Exercice 2 : Composition des ions} \hfill \textbf{/5 points}
	
	\vspace{3mm}
	
	L'atome d'aluminium possède 13 protons et 13 électrons. Il peut former l'ion aluminium \ch{Al^{3+}}.
	
	\vspace{3mm}
	
	\textbf{1.} L'ion aluminium est-il un cation ou un anion ? Justifier. \hfill \textit{(1 pt)}
	
	\cadreblanc{15}
	
	\textbf{2.} Combien de protons possède l'ion aluminium \ch{Al^{3+}} ? \hfill \textit{(1 pt)}
	
	\cadreblanc{12}
	
	\textbf{3.} Combien d'électrons possède l'ion aluminium \ch{Al^{3+}} ? \hfill \textit{(1,5 pts)}
	
	\cadreblanc{15}
	
	\textbf{4.} L'atome d'aluminium a-t-il gagné ou perdu des électrons pour devenir un ion ? Combien ? \hfill \textit{(1,5 pts)}
	
	\cadreblanc{15}
	
	
	
	\hrule
	\vspace{3mm}
	
	\textbf{Exercice 3 : Compléter le tableau} \hfill \textbf{/5 points}
	
	\vspace{3mm}
	
	Compléter le tableau suivant en calculant le nombre de protons et d'électrons pour chaque particule :
	
	\vspace{5mm}
	
	\begin{center}
		\begin{tabular}{|l|c|c|c|c|}
			\hline
			\textbf{Particule} & \textbf{Nombre de protons} & \textbf{Nombre d'électrons} & \textbf{Atome ou ion ?} & \textbf{Formule} \\
			\hline
			Calcium & 20 & 20 & & \\
			\hline
			& 20 & 18 & & \ch{Ca^{2+}} \\
			\hline
			Oxygène & 8 & 8 & & \\
			\hline
			& 8 & 10 & & \ch{O^{2-}} \\
			\hline
			& 29 & 27 & & \ch{Cu^{2+}} \\
			\hline
		\end{tabular}
	\end{center}
	
	\vspace{5mm}
	
	\hrule
	\vspace{3mm}
	
	
	\textbf{Exercice 4 : Tests de reconnaissance des ions} \hfill \textbf{/6 points}
	
	\vspace{3mm}
	
	On dispose de trois solutions A, B et C. On réalise des tests pour identifier les ions présents.
	
	\vspace{3mm}
	
	\textbf{Résultats obtenus :}
	
	\begin{itemize}
		\item \textbf{Solution A :} En ajoutant de l'hydroxyde de sodium, on observe un précipité rouille.
		\item \textbf{Solution B :} En ajoutant du nitrate d'argent, on observe un précipité blanc qui noircit à la lumière.
		\item \textbf{Solution C :} En ajoutant de l'hydroxyde de sodium, on observe un précipité bleu.
	\end{itemize}
	
	\vspace{3mm}
	
	\textbf{1.} Quel ion est présent dans la solution A ? Justifier. \hfill \textit{(2 pts)}
	
	\cadreblanc{20}
	
	\textbf{2.} Quel ion est présent dans la solution B ? Justifier. \hfill \textit{(2 pts)}
	
	\cadreblanc{20}
	
	\textbf{3.} Quel ion est présent dans la solution C ? Justifier. \hfill \textit{(2 pts)}
	
	\cadreblanc{20}
	
	
	
	\hrule
	\vspace{3mm}
	
	\textbf{Exercice 5 : Identifier les bons tests} \hfill \textbf{/4 points}
	
	\vspace{3mm}
	
	Pour chaque ion, indiquer le réactif à utiliser pour l'identifier et la couleur du précipité obtenu.
	
	\vspace{3mm}
	
	\textbf{1.} Ion fer (II) \ch{Fe^{2+}} \hfill \textit{(1 pt)}
	
	Réactif : \dotfill
	
	Couleur du précipité : \dotfill
	
	\vspace{5mm}
	
	\textbf{2.} Ion zinc \ch{Zn^{2+}} \hfill \textit{(1 pt)}
	
	Réactif : \dotfill
	
	Couleur du précipité : \dotfill
	
	\vspace{5mm}
	
	\textbf{3.} Ion chlorure \ch{Cl-} \hfill \textit{(1 pt)}
	
	Réactif : \dotfill
	
	Couleur du précipité : \dotfill
	
	\vspace{5mm}
	
	\textbf{4.} Ion cuivre (II) \ch{Cu^{2+}} \hfill \textit{(1 pt)}
	
	Réactif : \dotfill
	
	Couleur du précipité : \dotfill
	
	\vspace{5mm}
	
	\hrule
	\vspace{3mm}
	
	\textbf{Exercice 6 : Analyse d'une situation} \hfill \textbf{/4 points}
	
	\vspace{3mm}
	
	L'atome de chlore possède 17 protons et 17 électrons. Il peut former l'ion chlorure \ch{Cl-}.
	
	\vspace{3mm}
	
	\textbf{1.} Combien de protons et d'électrons possède l'ion chlorure \ch{Cl-} ? \hfill \textit{(2 pts)}
	
	\cadreblanc{10}
	
	\textbf{2.} L'ion chlorure est-il un cation ou un anion ? Justifier votre réponse. \hfill \textit{(2 pts)}
	
	\cadreblanc{10}
	
	\vfill
	
	\begin{center}
		\Large{\textbf{Barème total : /30 points}}
	\end{center}
	
\end{document}